\documentclass{article}
\usepackage{template/iclr2027_conference}
\usepackage{times}
\usepackage[T1]{fontenc}
\usepackage{amsmath,amssymb,graphicx,booktabs,longtable,array,calc}
\usepackage{hyperref}
\usepackage{etoolbox}
\hypersetup{colorlinks=true,linkcolor=blue,citecolor=blue,urlcolor=blue}
\providecommand{\tightlist}{\setlength{\itemsep}{0pt}\setlength{\parskip}{0pt}}
\makeatletter
\patchcmd{\@maketitle}{Under review as a conference paper at ICLR 2027}
  {Research draft --- not submitted}{\relax}{\errmessage{Draft header patch failed}}
\patchcmd{\@maketitle}{Anonymous authors\\Paper under double-blind review}
  {Research draft\\11 September 2026; not submitted for review}
  {\relax}{\errmessage{Draft author patch failed}}
\makeatother

\title{When Predictive Features Miss Motion: Measuring Temporal Accessibility in Skeleton JEPA}
\begin{document}
\maketitle
\begin{abstract}
Prediction of learned features need not preserve useful motion information. We connect two gait studies and implement a measurement path toward selective future distillation. In a retained skeleton-JEPA study with 625 clips from 93 videos, a matched movement readout falls from 0.223 \(R^2\) at initialization to 0.101--0.114 after training, despite consistent hidden-feature correspondence. In a repaired experiment with 50 clips from 43 videos, adding skeletons to an RGB reference yields -0.00024242 \(R^2\) on contextual video-teacher prediction. RGB complementarity, however, differs from skeleton-only accessibility. Our new source-held panel measures the predictive increment of a skeleton history block beyond current posture and observation support. Its primary effect is +0.001994 \(R^2\), with a conditional 95\% source-bootstrap interval of {[}-0.006673, +0.012608{]}: no supported temporal lead. The block also changes confidence features and regularization, so this estimate does not isolate motion. Exact reflection/time-reversal constructions and a controlled readout fixture separate geometric consistency from direction-sensitive information. We propose selecting future teacher directions for extra student-accessible temporal value, with an exact no-transfer option. The current contribution is a calibrated measurement framework and empirical motivation; improved student learning and independent confirmation remain to be demonstrated.
\end{abstract}
\section{Introduction}\label{introduction}

A predictive representation is useful relative to the distinctions it preserves. Action recognition may group walking examples despite changes in view or which leg moves faster. Measuring left--right movement requires those differences; anticipating a future position can require recent direction even when current posture is identical. One feature-prediction task need not preserve all three.

Joint-Embedding Predictive Architectures predict target representations from partial context. S-JEPA \citep{abdelfattah2024sjepa} applies this idea to hidden skeleton features, evaluated primarily through action recognition. Avoiding reconstruction of every noisy coordinate can be useful, but leaves open which movement quantities remain accessible.

We connect two development experiments. A laterality study finds improved hidden-feature correspondence but worse linear access to a signed speed contrast. Repaired future-innovation Experiment 0 instead finds no supported skeleton gain after an RGB reference when predicting contextual video-teacher features. Different targets, cohorts and \(R^2\) denominators prevent ranking the two results.

RGB complementarity also differs from skeleton-only accessibility. Shared motion information may be predictable from skeletons without improving a model that already sees RGB. Information requiring both modalities may improve their joint predictor while remaining inaccessible to a skeleton-only student. Distillation therefore needs evidence aligned with the student's inputs.

We implement a bounded source-held panel of a history-feature block beyond current posture and support, plus symmetry and temporal-observability calibrations. We then specify selection of future teacher directions as a subsequent experiment. The measured contribution is the separation and testing of these evidence levels; useful student transfer and broader world-modeling claims remain unestablished.

\section{Observables, representations and prediction boundaries}\label{observables-representations-and-prediction-boundaries}

\subsection{A movement measurement with known reflection behavior}\label{a-movement-measurement-with-known-reflection-behavior}

Both cohorts derive from GAVD \citep{ranjan2024gavd}, a video dataset for gait analysis. The laterality cohort contains 625 accepted clips from 93 source videos. For each of five bilateral landmark pairs, both sides must be observed at both ends of at least eight recorded transitions. Coordinates are centered at the pelvis and normalized by body width. Let \(m_{L,k}\) and \(m_{R,k}\) be median coordinate speeds, using the original time intervals. The scalar target is

\[
y=\frac{1}{5}\sum_{k=1}^{5}
\frac{m_{L,k}-m_{R,k}}{m_{L,k}+m_{R,k}+10^{-8}}.
\]

The pairs are shoulders, knees, ankles, heels and foot tips. Hips establish the origin and are excluded because centering equalizes their speed magnitudes. This target describes relative coordinate speed, not phase, contact, loading or clinical affected side. Image-normalized coordinates and inferred depth are not metric three-dimensional measurements.

An anatomical reflection M flips the centered horizontal coordinate and swaps left/right identities, confidence and validity. Then \(y(MH)=-y(H)\). Reversing a trajectory and its physical sampling intervals preserves speed magnitudes, so \(y(RH)=y(H)\). Laterality is therefore reflection-odd and time-even: it changes sign under one operation and remains unchanged under the other. Recovering it does not establish direction-sensitive forecasting.

\subsection{A contextual video-teacher target}\label{a-contextual-video-teacher-target}

Experiment 0 predicts a projected, 256-dimensional person-region representation from a frozen V-JEPA 2.1 teacher. Inputs cover frames 0--31; selected target tokens correspond to frames 38--39, but were encoded using the full 64-frame clip. The target can consequently reflect observations after its nominal frames. This is a declared contextual-feature task, not a measurement of an isolated future physical state. The nominal eight-frame horizon also corresponds to 0.267--0.800 seconds across the cached recordings.

The cache contains RGB/nuisance inputs, prefix skeletons, projected pooled targets and matching metadata. It does not contain mirrored teacher responses, dense token maps or future skeleton trajectories. New parity-resolved teacher features or explicitly bounded future-block targets require new encoding and corresponding provenance and time-boundary checks. V-JEPA 2.1's dense-feature design \citep{murlabadia2026vjepa21} motivates such targets but does not establish their gait-specific usefulness.

\subsection{Student-accessible temporal prediction}\label{student-accessible-temporal-prediction}

Let H contain skeleton history and permitted observation/timing channels. Define C as a fixed function of H retaining current posture and observation support. For a target \(Y_h\) at a declared horizon, write

\[
u_h(H)=\mathbb{E}[Y_h\mid H]-\mathbb{E}[Y_h\mid C].
\]

For finite second moments and a common squared-error metric, the population risk reduction equals \(\mathbb{E}\|u_h(H)\|^2\). This follows by decomposing \(Y_h-\mathbb{E}[Y_h\mid C]\) into the residual conditional on H and \(u_h\); their expected cross product is zero. These are standard conditional-expectation identities. They are not a new theorem, mutual-information estimator or guarantee for a regularized predictor trained on a small cohort.

The operational quantity is instead a difference between held-source errors of declared finite model families. It can be negative because of estimation error, regularization or an unsuitable representation. A null linear increment does not prove that the population conditional means coincide. Moreover, ordinary squared-error distillation already estimates a conditional mean in the ideal large-data limit. Selective targets must demonstrate an advantage in finite-data learning, robustness or downstream utility.

\begin{figure}[tbp]
\centering
\includegraphics[width=\linewidth]{figures/02_information_boundary.pdf}
\caption{Future observations supply training targets and scoring outcomes. They do not enter the evaluated student's prefix. A current-state reference and a history predictor must share the same declared support information.}
\label{fig:1}
\end{figure}

\section{A source-held accessibility panel}\label{a-source-held-accessibility-panel}

We implement a new development measurement using the existing 50 clips and 43 source videos. It preserves the teacher cache, source partitions, target units, raw skeleton controls, input-support preprocessing and scoring from direct-v3. It changes the question by replacing the RGB reference inputs with information available to a skeleton-only student.

The support panel uses 331 fixed inputs: per-joint confidence and validity means in four ordered eight-frame bins, frame-31 confidence and validity, and decoded frame rate. The posture panel adds the 66 valid x/y coordinates at frame 31, giving 397 inputs. Invalid endpoint coordinates are imputed using applicable training sources only. Coordinates inherit prefix-derived normalization; this reference is a prefix-normalized endpoint rather than an independently measured instantaneous image. Observation support deliberately remains available across the prefix because otherwise tracking patterns could be mistaken for a coordinate history benefit.

The added skeleton block contains 924 fixed features: ordered-bin coordinates, valid adjacent-frame velocities, confidence and observation/transition support. Velocities use only valid adjacent observations; imputation never creates movement. Raw histories form four arms before the common constructor: original, four-frame-block shuffle, different-source clip mismatch within the current partition, and no-skeleton coordinates/confidence with original validity. All arms receive the recipient's same reference block. Donor matching uses declared context metadata and no target values.

The common reference does not exhaust observation-quality information. Its confidence means include all frames, while S averages confidence over valid observations and the no-skeleton arm zeros those values. The primary contrast therefore concerns the entire declared coordinate/confidence feature block, including its additional regularized routes, rather than isolated motion.

Shuffle changes observation support and can create new block-boundary transitions; mismatch changes which body's observations are supplied. A favorable control comparison narrows explanations but does not isolate causal dynamics. Multiple noisy postures can also improve state estimation without requiring a learned direction-of-time mechanism.

Each fit minimizes a source-weighted, summed objective,

\[
\sum_i w_i\|Y_i-b-X_iW_x-S_iW_s\|^2
+\lambda_x\|W_x\|_F^2+\lambda_s\|W_s\|_F^2.
\]

Here X is the safely transformed reference, S the skeleton summary and Y the target standardized with current training-partition statistics. The squared Frobenius norm is the sum of squared matrix entries. Source weights give every video equal total influence and are normalized to sum to the training clip count. The intercept is unpenalized. The solve uses float64 arithmetic. Entirely unobserved, constant or insufficiently supported input columns have zero influence in all partitions; target standardization and masks remain separate.

Five outer source folds measure excluded-source prediction. Three inner source folds choose positive penalties from \(\{0.1,1,10,100,1000,10000\}\), using the same 36 joint combinations per arm. A separately selected shared reference is also an exact \texttt{baseline\_only} candidate. Inner losses pool source-weighted error sums consistently; ties within the frozen tolerance prefer the baseline, then stronger penalties. Failed required candidates prevent a complete scientific measurement. The model is deterministic and is not replicated under several seed labels.

The named primary contrast is real-minus-no-skeleton in the posture panel. The support panel and other arm comparisons are descriptive. A prospective development lead requires the primary paired 95\% interval above zero and positive point differences against shuffle and mismatch. This new exploratory rule does not revise direct-v3's original scientific thresholds or its STOP.

For teacher feature j, predictive \(R^2\) compares source-weighted squared prediction error with the error of the outer-training-mean reference. In training-standardized units that reference is zero. Valid target dimensions are the intersection of training-derived masks across outer folds; scores pool predictions before averaging featurewise \(R^2\). Two thousand paired bootstrap draws resample whole sources, retaining multiplicity and all associated clips. The intervals condition on the saved models and exclude repeated fitting, selection and adaptive redesign.

\section{Reflection and temporal observability as calibration}\label{reflection-and-temporal-observability-as-calibration}

Reflection and time reversal provide interpretable tests when their action on the observable is explicit. M and R commute under the declared landmark and timestamp transformations and each squares to identity. For any feature function h and signs \(a,b\in\{-1,+1\}\), define

\[
P_{a,b}h(H)=\frac{h(H)+a h(MH)+b h(RH)+ab h(MRH)}{4}.
\]

The projected feature gains factor a under M and b under R. This construction separates four transformation types, but its exactness is algebraic. Zero features satisfy the same rules. Feature energy and rank can detect that particular collapse, yet positive energy still does not establish useful movement information.

Our deterministic calibration checks the two involutions, their commutation, physical-interval reversal and four observable parity types. Group errors are zero on the constructed fixture; the largest recorded projector error is \(5.55\times10^{-17}\), below the prespecified \(10^{-12}\) tolerance. A zero-feature control has exact symmetry and zero predictive \(R^2\) on the balanced nonzero target. These results verify the construction rather than a learned model.

A second fixture has 48 generated sources, each contributing a closed path and its reversal. Both paths end at the same posture and have the same coordinate distribution and observation support, but their last signed vertical velocities are opposite. Sources 0--35 train a fixed-penalty linear readout; sources 36--47 are held out. The target is that last observed velocity, which defines a possible constant-velocity continuation. It is not measured future human motion.

{\def\LTcaptype{none} % do not increment counter
\begin{longtable}[]{@{}
  >{\raggedright\arraybackslash}p{(\linewidth - 2\tabcolsep) * \real{0.4286}}
  >{\raggedleft\arraybackslash}p{(\linewidth - 2\tabcolsep) * \real{0.5714}}@{}}
\toprule\noalign{}
\begin{minipage}[b]{\linewidth}\raggedright
Synthetic readout inputs
\end{minipage} & \begin{minipage}[b]{\linewidth}\raggedleft
Held-source \(R^2\)
\end{minipage} \\
\midrule\noalign{}
\endhead
\bottomrule\noalign{}
\endlastfoot
Current posture & 0.000000 \\
Coordinate mean, standard deviation, absolute changes and support & 0.000000 \\
Observation support only & 0.000000 \\
Ordered signed velocities & 0.999999999978 \\
\end{longtable}
}

The generator deliberately places opposite targets behind identical order-even summaries. This supplies a precise failure case for those summaries, not evidence that all temporal pooling is order-blind: contextual tokens may encode order before pooling. No JEPA training, GAVD forecasting or reflected-video encoding occurs in this calibration.

\section{Empirical evidence and its limits}\label{empirical-evidence-and-its-limits}

\subsection{Hidden-feature correspondence and movement readout disagree}\label{hidden-feature-correspondence-and-movement-readout-disagree}

The retained laterality grid contains five training conditions across five source folds and five seeds. With a common expanded readout, all trained teachers score below their matched initialization in every seed. The expanded summary includes means, standard deviations, absolute adjacent-block changes and support, giving 2,890 inputs compared with 960 for means alone.

{\def\LTcaptype{none} % do not increment counter
\begin{longtable}[]{@{}
  >{\raggedright\arraybackslash}p{(\linewidth - 6\tabcolsep) * \real{0.2143}}
  >{\raggedleft\arraybackslash}p{(\linewidth - 6\tabcolsep) * \real{0.2857}}
  >{\raggedleft\arraybackslash}p{(\linewidth - 6\tabcolsep) * \real{0.2857}}
  >{\raggedright\arraybackslash}p{(\linewidth - 6\tabcolsep) * \real{0.2143}}@{}}
\toprule\noalign{}
\begin{minipage}[b]{\linewidth}\raggedright
Laterality representation
\end{minipage} & \begin{minipage}[b]{\linewidth}\raggedleft
Mean \(R^2\)
\end{minipage} & \begin{minipage}[b]{\linewidth}\raggedleft
Trained minus initial
\end{minipage} & \begin{minipage}[b]{\linewidth}\raggedright
Retained 95\% source interval
\end{minipage} \\
\midrule\noalign{}
\endhead
\bottomrule\noalign{}
\endlastfoot
Matched initialization & 0.222544 & --- & --- \\
Motion random & 0.113725 & -0.108819 & {[}-0.170093, -0.041219{]} \\
MAMP-style mask & 0.112890 & -0.109653 & {[}-0.165071, -0.051351{]} \\
Motion mixture & 0.114209 & -0.108334 & {[}-0.170640, -0.042307{]} \\
Region random & 0.109446 & -0.113098 & {[}-0.166192, -0.059623{]} \\
Connected region & 0.100777 & -0.121766 & {[}-0.182787, -0.055281{]} \\
\end{longtable}
}

Meanwhile, correct-clip targets yield lower predictor error in all 375 trained diagnostic checks, compared with 33 of 75 distinct initial checks. This supports learning of correspondence, but neither identifies its mechanism nor turns repeated checks on the same models into independent samples. Mask comparisons against their respective random references have intervals crossing zero.

Several explanations remain open. Expanding the initial readout increases \(R^2\) by 0.151716, but simultaneously changes dimension, temporal summaries and support information. The largest tested ridge penalty wins for 49/125 teacher and 96/125 online expanded readouts, motivating a wider shared training-only grid. Input preparation also differs from target calculation: it fills gaps and resizes by sequence index, whereas the target retains original timing and gaps. The retained path comparison has 70.4\% sign agreement and \(R^2\) 0.218. This discrepancy neither bounds achievable accuracy nor identifies a single processing cause.

These are readout-access findings, not proof that motion information was erased. The current review recomputes retained seed and contrast arithmetic, but raw laterality checkpoints and prediction rows are absent from the copied bundle. The reported bootstrap intervals are retained results, not newly regenerated intervals. Original reflection experiments and the later mask grid also use different training/readout recipes and should not be combined as one factorial intervention.

\subsection{Repaired RGB complementarity remains unsupported}\label{repaired-rgb-complementarity-remains-unsupported}

The original residual implementation suffered from unsupported input directions, extreme scaling of newly observed missingness, and forced selection among corrections that all lost to the inner reference. Direct-v3 replaces that construction with the joint model, safe preprocessing and exact fallback. Its calibration and independent prediction reconstruction support interpreting the repaired negative result.

{\def\LTcaptype{none} % do not increment counter
\begin{longtable}[]{@{}
  >{\raggedright\arraybackslash}p{(\linewidth - 4\tabcolsep) * \real{0.2727}}
  >{\raggedleft\arraybackslash}p{(\linewidth - 4\tabcolsep) * \real{0.3636}}
  >{\raggedleft\arraybackslash}p{(\linewidth - 4\tabcolsep) * \real{0.3636}}@{}}
\toprule\noalign{}
\begin{minipage}[b]{\linewidth}\raggedright
Direct-v3 predictor
\end{minipage} & \begin{minipage}[b]{\linewidth}\raggedleft
Predictive \(R^2\)
\end{minipage} & \begin{minipage}[b]{\linewidth}\raggedleft
Gain over shared RGB reference
\end{minipage} \\
\midrule\noalign{}
\endhead
\bottomrule\noalign{}
\endlastfoot
RGB/nuisance reference & 0.334215168 & --- \\
Real skeleton & 0.333972751 & -0.000242417 \\
Time shuffle & 0.334246668 & +0.000031500 \\
Clip mismatch & 0.335452594 & +0.001237425 \\
No skeleton, validity retained & 0.334215168 & 0.000000000 \\
\end{longtable}
}

The matched increment is -0.000242417, with conditional 95\% interval {[}-0.001474019, +0.000919797{]} and 36.35\% positive paired draws. Fourteen of twenty selected fold-arm models use the exact baseline. All 740 pooled candidates are valid. All selected joint models retain the reference's RGB penalty of 100, so a changed selected RGB penalty does not explain these particular increments.

The original direct-v2 scores remain the result of the earlier implementation. Its post-hoc weight removal and zero-initialization interventions are diagnostic, not independent evidence of a repaired improvement. The increase in the RGB reference between protocols is likewise not a skeleton gain.

Unlike the teacher-feature experiment, laterality's scalar \(R^2\) uses the weighted held-out mean in its denominator. We preserve both definitions and do not compare their absolute levels. Both cohorts are development data, and source holdout does not establish participant holdout across recordings.

\subsection{Student-accessibility measurement}\label{student-accessibility-measurement}

The new panel completes all required fits and retains all 256 target dimensions. All 1,480 pooled candidates are valid; 17 of 40 selected fold-arm models use exact baseline fallback. Read-only verification reloads and refits selected linear models and preprocessing, reconstructs 102,400 held-out prediction rows and 16,000 bootstrap rows, and checks the saved report. Parent and source-run artifacts remain unchanged. Inner candidate losses are checked for pooling and selection, but are not independently refitted in this verification.

{\def\LTcaptype{none} % do not increment counter
\begin{longtable}[]{@{}lrr@{}}
\toprule\noalign{}
Cached-panel predictor & Support reference & Posture reference \\
\midrule\noalign{}
\endhead
\bottomrule\noalign{}
\endlastfoot
Shared reference & 0.008229092 & 0.005890997 \\
Real skeleton & 0.013864827 & 0.013261709 \\
Time shuffle & 0.012433535 & 0.007692185 \\
Clip mismatch & 0.008229092 & 0.005890997 \\
No skeleton, validity retained & 0.007295537 & 0.011267966 \\
\end{longtable}
}

\emph{Table 4. Predictive \(R^2\) on the existing contextual teacher target. The posture panel's real-minus-no-skeleton contrast was named primary before fitting.}

The primary estimate is \textbf{+0.001993744 \(R^2\)}, with paired 95\% interval \textbf{{[}-0.006672679, +0.012608431{]}} and \textbf{68.35\%} positive draws. The result is \texttt{no\_supported\_temporal\_lead} under the prospective exploratory rule. Real-minus- baseline is +0.007370712, but the real model selects stronger reference-block regularization in two folds, and the validity-retaining no-skeleton model gains +0.005376968 itself. The reference difference therefore cannot be attributed entirely to coordinate history. The matched increment is the relevant declared comparison, and it remains uncertain.

Adversarial reconstruction identified an additional interpretation limit: 182 finite confidence summaries differ between X's raw means and S's valid-conditioned means. Real-arm fits retain 132 confidence columns, versus zero in no-skeleton. Thirty-one confidence columns exactly duplicate standardized X columns in every fold. Duplicate columns with penalties \(\lambda_x\) and \(\lambda_s\) have effective penalty \((\lambda_x^{-1}+\lambda_s^{-1})^{-1}\), so they alter regularization as well as nominal representation. The current evidence does not quantify their contribution to the estimate. The preserved negative conclusion does not establish that confidence and recording quality were fully controlled.

The secondary support panel gives real-minus-no-skeleton +0.006569290, interval {[}-0.000071482, +0.014709426{]}, with 97.3\% positive draws. Its real-minus-shuffle contrast is only +0.001431291, interval {[}-0.003920839, +0.008541424{]}. It also omits current posture from the reference. Selecting this more attractive result after fitting would change the scientific question. Both panels show that these fixed skeleton summaries predict little of the full contextual teacher vector under this training regime; they leave target suitability, model family and data quantity unresolved.

\begin{figure}[tbp]
\centering
\includegraphics[width=\linewidth]{figures/06_cached_panel_results.pdf}
\caption{Differences in predictive \(R^2\) with paired 95\% source-bootstrap intervals, conditional on fitted models. Teal marks the prospectively designated posture-panel comparison. The secondary panel does not replace it; every displayed interval includes zero.}
\label{fig:2}
\end{figure}

\section{A prospective selective-distillation method}\label{a-prospective-selective-distillation-method}

Before isolating coordinate motion in a new panel, include every confidence, validity and transition-support summary once in a common block, identically across arms. Vary only coordinates and displacements; use that common block as the confidence/support-only reference and control its coefficients and penalty. These are prospective amendments requiring a new frozen experiment, not changes to the completed comparison.

The next method experiment would select teacher directions by their source-held temporal value. Let U contain a small number of constrained teacher directions, learned using training sources only. Cross-fitted current-state predictions \(\widehat m_C\) define residual targets in common raw teacher units. A student would minimize an auxiliary objective of the form

\[
\mathcal L=\mathcal L_{\text{S-JEPA}}+
\beta\,\mathbb E\big\|g_\theta(H)-
U^\top[Y_h-\widehat m_C(C)]\big\|^2.
\]

This is a proposed experiment, not a completed student implementation or result. Ranks and weights include an exact zero-transfer choice. Normalization, cross-fitting, direction learning and selection must stay inside the relevant training partitions. U must have a fixed scale, such as orthonormal columns, to prevent trivial loss reduction by shrinking targets. The comparison should retain both original-target measurements and a common downstream outcome.

Reflection-resolved targets are one candidate, not a requirement that all useful motion be asymmetric. A teacher's even/odd components require actual paired encodings of original and reflected video in a compatible feature basis. They cannot be inferred by negating an existing pooled vector. Their motion content must be tested because teacher reflection sensitivity can encode view or appearance.

A focused target-construction study could pair similar current postures with different incoming velocities. Within-recording future-feature differences may attenuate shared appearance, but this requires more same-source windows and an empirical cancellation check. Pair construction must use prefix information only, with sources split before pairing. Controlled motion under independently varied appearance provides a useful mechanism test; naturally matched videos remain observational evidence.

The decisive student comparison holds architecture, data and training opportunity fixed across ordinary S-JEPA, direct motion prediction, complete teacher distillation, unselected residual distillation, a dimension-matched target chosen without temporal information, and the proposed selection. Evaluation must include observable future motion, current-state and simple velocity baselines, and independent sources or participants. A lower selected-code loss alone cannot establish useful transfer.

\section{Related work and contribution boundary}\label{related-work-and-contribution-boundary}

Motion-sensitive skeleton pretraining is established by MAMP \citep{mao2023mamp}; SLiM \citep{do2026slim} studies compact feature prediction and semantic tube masks. MC-JEPA \citep{bardes2023mcjepa} combines motion and content learning. Our local results do not reproduce or invalidate these complete systems; they motivate evaluating specific observables alongside the learned prediction objective.

Invariant/equivariant separation also has clear precedents in SIE \citep{garrido2023sie} and seq-JEPA \citep{ghaemi2025seqjepa}. The four-parity construction is a standard finite-group projection, used here as an interpretable calibration. Neither it nor an added reflection loss constitutes the proposed novelty.

The Modality Focusing Hypothesis \citep{xue2023modality} studies when crossmodal knowledge transfer helps; C2VL \citep{chen2024c2vl} transfers visual-language information into skeleton representations. Future privileged information is central to Overlooked Poses \citep{sun2022overlooked} and \href{https://openreview.net/pdf?id=P6F4MxtOKp}{Spectral-Guided Physical Dynamics Distillation}, the latter including spectral enhancement and human-motion evaluation. The prospective distinction is therefore narrow: selection by incremental temporal accessibility, explicit observable checks and demonstrated avoidance of unsuitable transfer. Conditional future representations themselves have a long history, including predictive-state regression \citep{hefny2015predictive}.

BioGait-VLM \citep{chen2026biogait} combines observed video, language and biomechanical information for gait assessment. Its temporal evidence aggregation is a further reason to avoid a broad gait-distillation novelty claim; our proposed future-prefix student comparison addresses a narrower question.

\section{Discussion}\label{discussion}

Hidden-feature correspondence, anatomical symmetry, temporal accessibility and student utility are distinct properties. The completed comparisons do not establish the last of these. The next decisive evidence is a source-held student comparison with bounded future targets, common motion outcomes and matched alternatives, followed by confirmation with participant identities or reliable geometry. The inspected 43-source cache refines that question without establishing broad transfer. Separate-horizon prediction would still not demonstrate recursive rollout, and these observations contain no actions or closed-loop decisions. The immediate objective is to test when transferring future information improves measured movement prediction, including a valid choice to transfer nothing.

\label{main-text-end}
\clearpage
\bibliographystyle{template/iclr2027_conference}
\bibliography{references}
\section*{Reproducibility and evidence statement}\label{reproducibility-and-evidence-statement}

The repository retains the original laterality and future-innovation evidence, source inventories, calibration definitions, separately versioned panel code, selected-model artifacts and numerical verification. Teacher-cache reuse is read-only. The current laterality check covers retained aggregate arithmetic; the repaired future-innovation check reconstructs predictions and paired scores. The accompanying tutorial and panel protocol document those boundaries and execution commands. The prospective selective student experiment has no reported real-data performance.

\section*{Ethics and AI-use statement}\label{ethics-and-ai-use-statement}

Video-derived poses can reveal personal information and retain demographic or recording biases. The coordinate-speed contrast has no demonstrated clinical diagnostic validity, and public annotation access does not authorize video redistribution. Source grouping does not guarantee distinct participants. AI assistance supported evidence organization, implementation, literature checking, figure programming and manuscript drafting. Authors remain responsible for source verification, experimental interpretation and the submitted claims.

This draft still requires complete human author verification before submission.

\end{document}
